% 本文件由 LaTeX 排版 Agent 根据有效 translation.json 自动生成。
% author=John Chrysostom; work_id=1910; title=讲道：论父啊，若是可能…

\chapter{讲道：论\Quote{父啊，若是可能…}}

% structure: p0001 source=h1 target=omitted reason=page-title-already-rendered-by-parent

\Emph{驳斥马尔西翁派与摩尼教}\par

论经文\Quote{父啊，若是可能，请让这杯离开我；但不要照我的意愿，而要照你的意愿：}并驳斥马尔西翁派与摩尼教；以及我们不应莽撞冒险，而应将天主的旨意置于一切意愿之上。\par

我最近给那些贪得无厌、损人利己的人猛击了一掌；这样做并非为了伤害他们，而是为了纠正他们；并非因为我憎恨这些人，而是因为我厌恶他们的恶行。因为医生也是如此，他切开脓肿，并非攻击受苦的身体，而是为了对付疾病和创伤。今天让我们稍作休息，使他们从痛苦中恢复过来，以免因持续受折磨而抗拒治疗。医生们也是这样行事的；切开之后，他们敷上膏药和药物，并过几天才设法止痛。效法他们的榜样，让我今天也为他们设计一些能从我的讲道中获益的办法，提出一个有关教义的问题，将我的言论集中到所宣读的经文上。因为我想，有许多人对基督说这些话的原因感到困惑；而且很可能，在场的异端者也会抓住这些话，从而动摇许多心性单纯的弟兄们。\par

为要筑起一道墙来抵挡他们的攻击，并使困惑者摆脱迷茫和混乱，让我们拿起所引用的经文，仔细思考这段经文，深入其含义的深处。因为仅仅阅读是不够的，除非加上知识。正如甘达刻的太监阅读，但直到有人来指示他所读内容的含义，他才从中获得很大的益处。因此，为使你们不至于同样，请留意所说的，运用你们的理智，让我看到你们的心思从其他思绪中解脱出来，让你们的眼睛敏锐，意向诚恳；让你们的灵魂从世俗忧虑中解放出来，免得我们把种子撒在荆棘上，或石头上，或路旁，而是要深耕肥沃的田地，从而收获丰盛的庄稼。因为如果你们如此留意所说的，你们将使我的劳作变得轻松，并有助于发现你们所寻求的。\par

那么，所宣读的经文\Quote{父啊，若是可能，请让这杯离开我？}是什么意思？这话是什么意思？因为我们应当先清楚解释这些话语的含义，然后才开启经文。那么这话是什么意思？\Quote{父啊，若是可能，请除去十字架。}你怎么说？祂不知道这事是否可能吗？谁敢这样说？但这话是出于无知者的口；因为加上\Quote{若是}这个词表明怀疑；但正如我说的，我们不应只留意话语，而应留意其含义，了解说话者的意图、原因和场合，并将这些综合起来，揭示隐藏的意思。既然如此，那不可言传的智慧，认识父如同父认识子，祂如何会不知道这事？因为关于祂受难的知识并不比关于祂本质的知识更大，而后者只有祂自己准确知道。因为祂说：\BibleQuote{就如父认识我，我也认识父}\BibleRef{若 10:15}。我为何只提到天主的独生子？因为就连先知们似乎也没有不知道这事，反而清楚地知道，并充满信心地预言这事必须成就，且必然发生。\par

请至少听一听，众人如何以各种方式宣告十字架。首先，圣祖雅各伯：当他将话语指向祂时，说道：\Quote{你从嫩枝中升起：}用\Quote{枝}字表示童贞女和玛利亚无玷的本性。接着，为指明十字架，他说：\BibleQuote{你卧下沉睡如狮子，又如幼狮；谁能唤醒祂？}\BibleRef{创 49:9}在此，他把死亡称为沉睡和睡眠，并在死亡中结合了复活，当他说：\Quote{谁能唤醒祂？}确实，除了祂自己，无人能唤醒——因此基督也说：\BibleQuote{我有权舍掉我的生命，也有权再取回来}\BibleRef{若 10:18}，又说：\BibleQuote{拆毁这座圣殿，三天之内我要把它重建起来。}\BibleRef{若 2:19}而\Quote{你卧下沉睡如狮子}这话是何意？因为狮子不仅清醒时可怕，连睡觉时也可怕；同样，基督不仅生前可怕，就是在十字架上和死亡的那一刻，也是可怕的，并在那时行了伟大的奇迹：使日光回转，劈开岩石，震动大地，撕裂帐幔，惊扰比拉多的妻子，定犹达斯的罪，因为他当时说：\BibleQuote{我出卖了无辜的血，犯了罪了。}\BibleRef{玛 27:4}比拉多的妻子也宣称：\BibleQuote{不要干涉那义人，因为我今天在梦中因祂受了许多苦。}\BibleRef{玛 27:19}黑暗笼罩大地，正午出现黑夜，然后死亡被消灭，它的暴政被摧毁：许多长眠的圣人的身体复活了。圣祖预先宣告这些事，并表明基督即使被钉十字架，仍是可怕的，于是说：\Quote{你卧下沉睡如狮子。}他没有说\Quote{你将沉睡}，而是说\Quote{你沉睡了}，因为这必然会发生。因为先知们常在许多地方把未来的事当作已过的事来预言。正如已经发生的事不可能不发生一样，尽管这事是未来的，它也不可能不发生。为此，他们以过去时的形式预言未来，以此表明它们不可能失败，必然应验。达味也是这样说话，指明十字架：\Quote{他们刺穿了我的手和我的脚。}他没有说\Quote{他们将刺穿}，而是说\Quote{他们刺穿了}\Quote{他们数算了我所有的骨头。}他不仅说了这些，还描述了士兵们所做的事：\Quote{他们彼此分了我的衣服，为我的外衣拈阄。}不仅如此，他还提到他们给祂吃苦胆，喝醋；因为他说：\Quote{他们给了我苦胆作食物，我渴时，他们给了我醋喝。}又有另一位说，他们用长枪刺了祂，因为\BibleQuote{他们要仰望他们所刺透的那一位。}\BibleRef{匝 12:10}依撒意亚又以另一种方式预言十字架说：\BibleQuote{祂被牵去宰杀，如同羔羊；又像母羊在剪毛者前缄默，祂也不开口。在祂受辱时，审判被夺去了。}\BibleRef{依 53:7}\par

2. 如今我恳请你们注意，这些作者每一位都如何好似谈论已过之事般发言，以此动词时态表明事件绝对无可避免的确凿性。同样，达味在描述这一审判座时也说：\Quote{外邦为什么怒吼？万民为什么妄想？地上的君王起来，首领聚集一起，反对主，反对祂的基督。}他不只提到审判、十字架和十字架上的事，也提到出卖祂的人，指明那是祂的亲密同伴和宾客。\Quote{因为，}他说，\Quote{那与我同食的人，竟抬脚踢我。}同样，他也预言了基督要在十字架上发出的呼声，说：\Quote{我的天主，我的天主，祢为什么舍弃了我？}也描述了埋葬：\Quote{祢把我放在极深的坑里，在黑暗处，在死荫里。}以及复活：\Quote{祢必不将我的灵魂遗弃在阴间，也不让祢的圣者见到腐朽。}和升天：\Quote{天主欢呼上升，主在号角声中上升。}以及坐在右边：\Quote{主对我主说：祢坐在我右边，等我使祢的仇敌作祢的脚凳。}但依撒意亚也宣告了原因，说：\Quote{祂因我百姓的过犯被带到死地，}\BibleRef{依 53:8}又因为众人如羊走迷，所以祂被牺牲。\BibleRef{依 53:6}接着他也提到结果，说：\Quote{因祂的创伤我们全都得了痊愈，}\BibleRef{依 53:5}以及\Quote{祂担当了许多人的罪。}\BibleRef{依 53:12}先知们既知道十字架、十字架的原因及其效果，也知道埋葬、复活、升天、出卖和审判，并且准确描述了一切。那么，那位派遣他们并命令他们说这些话的，反而对这些事一无所知吗？哪个有理智的人会这样说？你们看见，我们不可只留意字句吗？因为令人困惑的经文不只是这一处，后面的更为费解。因为祂说什么？\Quote{父啊！若是可能，就让这杯离开我吧。}在此，祂显得不仅好像无知，而且好像不愿接受十字架：因为祂说这话。\Quote{若是许可，求祢不要让我受十字架和死亡。}然而，当宗徒之长伯多禄对祂说：\Quote{主，千万不要这样，这事绝不会临到祢身上，}祂严厉地斥责他说：\Quote{撒殚，退到我后面去！你是我的绊脚石，因为你所体会的不是天主的事，而是人的事。}\BibleRef{玛 16:22}虽然不久之前祂刚称他有福。逃避十字架在祂眼中是如此可怕，以致那位从父领受了启示、祂曾称有福、并领受天国之钥匙的人，祂竟称他为撒殚、绊脚石，并指责他不体会天主的事，因为他对祂说：\Quote{主，千万不要这样，这事绝不会临到祢身上}——指十字架。那么，祂既然这样痛斥门徒，对他大发雷霆，甚至称他为撒殚（在赐予他如此大赞之后），只因他说了\Quote{避免十字架}，祂又怎能不愿被钉呢？此后，当祂描绘善牧的比喻时，又怎能宣称这是祂德行的特别证明，即祂应为羊群牺牲，如此说：\Quote{我是善牧：善牧为羊舍掉自己的性命？}\BibleRef{若 10:11}祂甚至没有停在那里，还加上说：\Quote{但雇工，不是牧人，看见狼来，就撇下羊逃跑了。}\BibleRef{若 10:12}如果牺牲自己是善牧的记号，而不愿如此行是雇工的记号，那么那自称为善牧的又如何恳求不要被牺牲呢？祂又怎能说：\Quote{我舍掉我的性命，是出于我自己？}因为如果你自己舍掉性命，又怎能恳求别人让你不捨呢？保禄又怎会因这一宣言而惊叹说：\Quote{祂虽具有天主的形体，却没有以自己与天主同等为持守的宝物，反而空虚自己，取了奴仆的形体，成了人的模样；且以人的形态出现，祂贬抑自己，听命至死，且死在十字架上。}\BibleRef{斐 2:6}而且祂自己又这样说：\Quote{为此父爱我，因为我舍掉我的性命，为再取回它。}\BibleRef{若 10:17}如果祂不愿舍弃性命，反而拒绝此事，恳求父，那么祂怎会因此被爱呢？因为爱是同心合意之人的事。保禄又怎么说：\Quote{彼此相爱，如同基督爱了我们，为我们舍弃了自己？}\BibleRef{弗 5:2}基督自己即将被钉时，说：\Quote{父啊，时候到了，请祢光荣祢的子，}\BibleRef{若 17:1}将十字架称为光荣：那么祂又怎能在那里催促十字架，而在这里拒绝它呢？因为十字架是光荣，请听福音作者所说：\Quote{圣神还没有赐下，因为耶稣还没有受到光荣。}\BibleRef{若 7:39}这句话的意思是：\Quote{恩宠还没有赐下，因为天主对人的仇恨还没有被摧毁，因为十字架尚未完成它的工作。}因为十字架摧毁了天主对人的仇恨，带来了和好，使地成为天，将人联合于天使，推倒了死亡的堡垒，瓦解了魔鬼的力量，熄灭了罪恶的权势，将世界从错误中解救出来，带回了真理，驱逐了恶魔，毁灭了庙宇，推翻了祭台，压制了祭献，培植了德行，建立了教会。十字架是父的旨意、子的光荣、圣神的喜乐、保禄的夸耀，\Quote{因为，}他说，\Quote{我绝不夸耀，除了夸耀我们主耶稣基督的十字架。}\BibleRef{迦 6:14}十字架比太阳更明亮，比光线更灿烂：因为当太阳昏暗时，十字架却光辉灿烂：太阳昏暗并非是因它熄灭，而是被十字架的光辉所胜。十字架折断了我们的债券，使死亡的监狱失效，它是天主之爱的明证。\Quote{天主竟这样爱了世界，甚至赐下自己的独生子，使凡信祂的人不至丧亡。}\BibleRef{若 3:16}保禄又说：\Quote{如果我们还在为仇敌的时候，因祂圣子的死，得以与天主和好。}\BibleRef{罗 5:10}十字架是不可攻破的城墙、不可损伤的盾牌、富人的保障、穷人的依靠、遭暗算者的防卫、受攻击者的武装、克制情欲的方法、获得德行的途径、奇妙惊人的记号。\Quote{这一世代寻求征兆；除了约纳的征兆外，再没有征兆赐给它。}\BibleRef{玛 12:39}保禄又说：\Quote{犹太人要求征兆，希腊人寻求智慧，但我们却宣讲被钉的基督。}\BibleRef{格前 1:22}十字架敞开了乐园，领进了强盗，将即将灭亡、连地都不配的人带进了天国。如此众多且继续由十字架产生的恩惠，而祂竟不愿被钉，我问你们？谁敢这样说？如果祂不愿，谁强迫了祂，谁迫使了祂？如果祂不愿承受，又何必事先派遣先知宣告祂将被钉呢？如果祂不愿被钉，又为何称十字架为杯？因为那是表达对某事渴望之人的话。正如杯对口渴者是甘甜的，同样，十字架对祂也是如此：因此祂也说：\Quote{我渴望又渴望，在受难以前同你们吃这一次逾越节晚餐，}\BibleRef{路 22:15}而祂这样说不是绝对意义上，而是相对而言，因为那晚之后，十字架正等着祂。\par

3. 那么，那称这事为光荣，并斥责门徒因为他试图阻止祂，且证明好牧人的要素在于他为羊群舍弃自己，并宣称他切望此事，甘愿迎向它的人，祂又怎会祈求它不要成就呢？如果祂不愿意，阻止那些为此而来的人有什么困难呢？但事实上，你看见祂急忙走向这事。至少当他们来到祂跟前时，祂说：\BibleQuote{你们找谁？}他们回答说：\BibleQuote{耶稣。}然后祂对他们说：\BibleQuote{我就是}，他们便退后，跌倒在地上。\BibleRef{若 18:6}这样，祂先使他们瘫痪，证明祂能逃脱他们的手，然后祂才交出自己，使你明白，祂不是被迫或勉强，或由于攻击者暴虐的力量而不情愿地屈服于此，而是甘愿地、有目的有渴望地，早已为此预备了多时。因此，先知们先被派遣，圣祖们预言了这些事，十字架通过言语和行为被预表。因为依撒格的祭献也向我们预示了十字架；为此基督也说：\BibleQuote{你们的父亲亚巴郎曾欢欣地看到我的光荣，他看见了，就喜乐。}\BibleRef{若 8:56}那时，圣祖看见十字架的图像就喜乐，而祂自己却反对它吗？同样，梅瑟在展示十字架形象时，战胜了阿玛肋克；人们可以在旧约中看到无数预先描绘十字架的事。那么，如果将被钉十字架的那位不愿意这事发生，这又是为什么呢？此后的话语更加令人困惑。因为祂说了\Quote{愿这杯从我这里过去}后，又加上：\BibleQuote{然而不要照我的意思，只要照你的意思。}\BibleRef{玛 26:39}因为从实际表达来看，我们发现两个意愿彼此对立：至少圣父愿意祂被钉十字架，而祂自己却不情愿。然而我们处处看到祂愿意并计划与圣父相同的事。因为当祂说：\BibleQuote{使他们合而为一，如同我们原为一体，叫他们也在我们内合而为一，}\BibleRef{若 17:11}这等于说圣父和圣子的旨意是一。当祂说：\BibleQuote{我对你们所说的话，不是凭自己说的，而是住在我内的圣父作祂自己的事，}\BibleRef{若 14:10}祂也表明同样的事。当祂说：\BibleQuote{我不是凭我自己来的}\BibleRef{若 7:28}和\BibleQuote{我凭自己不能做什么}\BibleRef{若 5:30}，祂这样说并非表明祂被剥夺了说话或行事的权柄（切勿这样想！），而是愿意证明祂的旨意，无论在言语、行为还是各类交往中，都与圣父完全相同，正如我多次证明过的。因为\BibleQuote{我不凭自己说}这句话不是权柄的废除，而是同意的证明。那祂在这里怎么说\BibleQuote{然而不要照我的意思，只要照你的意思}呢？也许我在你们心中引起了很大的冲突，但请保持警觉：虽然我说了许多话，但我知道你们的热情仍然新鲜；因为论述现在正急于给出解答。那么为什么使用这种说话方式呢？请留心听：降生成人的道理很难接受。因为祂慈爱的超然尺度与屈尊就卑的伟大充满了敬畏，需要很多预备才能被接受。请思考，听到并学习如下之事是何等伟大：天主，那不可言说、不可朽坏、不可领悟、不可眼见、不可理解的，大地四极在祂手中，祂观看大地，大地便战栗，祂触摸群山，群山便冒烟，连革鲁宾也不能承受祂屈尊就卑的重量，而以翅膀的庇护掩面；这位超越所有理解、难倒一切计算的天主，越过天使、总领天使和所有天上的神体，屈尊成为人，取用了由土和泥形成的肉身，进入童贞女的子宫，在那里被怀胎九个月，以奶喂养，并承受人所应受的一切。既然将要发生的事如此奇特，以至于许多人在它发生后仍不相信，祂首先派遣先知们预先宣布这件事。例如，圣祖预言说：\BibleQuote{我儿，你从嫩芽中生出；你躺卧如公狮，沉睡如母狮}\BibleRef{创 49:9}；依撒意亚说：\BibleQuote{看，童贞女将怀孕生子，人要称祂的名叫厄玛奴耳}\BibleRef{依 7:14}；又在别处说：\BibleQuote{我们看见祂像幼儿，像干地上的根}\BibleRef{依 53:2}；祂用干地意指童贞女的子宫。又说：\BibleQuote{有一个婴孩为我们诞生了，有一个儿子赐给了我们}\BibleRef{依 9:6}；又说：\BibleQuote{从叶瑟的根茎将生出一根枝条，从他的根上将开出一朵花}\BibleRef{依 11:1}。巴路克在\WorkTitle{耶肋米亚书}中说：\BibleQuote{这就是我们的天主，没有别的可与之相比；祂发现了知识的一切道路，将智慧赐给祂的仆人雅各伯，和祂所爱的以色列。此后，祂在地上出现，与人往来。}达味预示祂的降生成人，说：\BibleQuote{祂如雨降落在羊毛上，如滴露落下大地}，因为祂无声无息、温和地进入了童贞女的子宫。\par

4. 但这些证据仍不足以证明；甚至当祂来临之后，为免所发生的事被视为幻象，祂不仅借视觉，更借时间的延续和经历人所固有的一切阶段来证实这事实。因为祂并非一次进入一个成熟且完全发育的人，而是进入童贞玛利亚的胎中，以便经历孕育、诞生、哺乳和成长的过程，并借时间的长度和成长阶段的多样性来确保所发生之事。甚至这些证据也并未结束；当祂带着血肉之体时，祂容许它体验人性的软弱：饥饿、口渴、睡眠和疲劳；最后，当祂走向十字架时，祂容许它经受肉身的痛苦。为此，汗珠从祂身上流下，有天使被发现来增强祂的力量，祂忧伤沮丧；因为在说出这些话之前，祂说：\BibleQuote{我的心灵忧闷，且极其忧伤，至死为止。}\BibleRef{玛 26:38} 若在这些事发生之后，魔鬼借着\Place{本都}的\Person{马尔基昂}、\Person{瓦伦提努斯}、\Place{波斯}的\Person{摩尼}和许多其他异端者的邪恶之口，妄图推翻降生成人的教义，并发出魔鬼般的言论，声称祂并未成为血肉，也未穿上血肉，而仅仅是幻想、幻象、演戏和伪装——尽管痛苦、死亡、埋葬、口渴都在大声反对这种教导——那么，假若这些事都未发生，魔鬼岂不是更广泛地散布这些邪恶的不敬教义吗？为此，正如祂饥饿、睡觉、疲劳、吃喝一样，祂也祈求免于死亡，从而彰显祂的人性，以及那不愿痛苦地脱离此生的人性软弱。因为若祂未说出任何这些事，就可能会有人说，若祂是人，祂理当体验人的情感。这些情感是什么呢？对于一个即将被钉十字架的人来说，就是恐惧、痛苦和脱离此生的痛苦：因为对周围事物魅力的感觉已植根于人性中；为此，为证明血肉之衣的真实性，并确认降生成人，祂以充分的证据显明人实际的情感。\par

这是一方面的考虑，但还有另一同等重要的方面。是什么呢？基督来到世上，愿意教导人一切德行：而导师不仅用言语，也用行动来教导；因为这是教师最好的教学方法。例如，舵手让学徒坐在身边，向他示范如何掌舵，同时也以行动配合言语，不单靠言语或只靠榜样；同样，建筑师让有意学习砌墙的人站在身旁，不仅以口授教导，也以行动示范；织工、绣匠、金匠、铜匠也是如此——每一门技艺的教师都兼用口头和实践来指导。既然基督亲自来教导我们一切德行，祂便既告诉我们当做之事，也亲自实行。祂说：\BibleQuote{因为，谁若实行并教导同样的，在天国里将被称为大的。}\BibleRef{玛 5:19} 现在请留意：祂命令人要谦卑和温和，并以言语教导了这一点；但请看祂如何也用行动来教导。祂说了\BibleQuote{神贫的人是有福的，温良的人是有福的}\BibleRef{玛 5:3} 之后，就显示这些德行当如何实践。祂是如何教导的呢？祂拿了一条毛巾，束上腰，洗了门徒的脚。\BibleRef{若 13:4} 有什么能比得上这谦卑呢？因为祂不再只用言语，也用行动来教导这德行。祂又以行动教导温和与忍耐。如何呢？祂被大司祭的仆人打脸，却说：\BibleQuote{我若说得不对，你指证那里不对；若说得对，你为什么打我？}\BibleRef{若 18:23} 祂命令人要为自己的仇敌祈祷：这也再次以行动来教导：当祂被钉在十字架上时，祂说：\BibleQuote{圣父，宽赦他们吧！因为他们不知道他们做的是什么。}\BibleRef{路 23:34} 因此，正如祂命令人祈祷，祂自己也祈祷，用祂自己不断的祈祷教导你也这样做。祂又命令我们善待恨我们的人，以德报怨：\BibleRef{玛 5:44} 祂也以自己的行动行了这事：因为祂从犹太人身上驱逐魔鬼，而那些人却说祂自己被魔鬼附身；祂恩待迫害自己的人，喂养那些图谋反对祂的人，带领那些渴望钉死祂的人进入祂的王国。祂又对门徒说：\BibleQuote{你们不要在腰包里带金、银、铜钱}\BibleRef{玛 10:9}，以此训练他们神贫；祂也以身作则教导这一点，说：\BibleQuote{狐狸有穴，天上的飞鸟有巢，但是人子却没有枕头的地方。}\BibleRef{玛 8:20} 祂没有桌子、住所或任何这类东西：不是因祂无法得到，而是为教导人走上那条路。同样地，祂也教导他们祈祷。他们对祂说：\BibleQuote{请教给我们祈祷。}\BibleRef{路 11:1} 因此祂也祈祷，使他们学习祈祷。但不仅需要学习祈祷，还需知道如何祈祷：为此祂传授了这样一篇祈祷：\BibleQuote{我们在天的圣父！愿你的名被尊为圣，愿你的国来临，愿你的旨意承行于地，如在天上一样。求你今天赏给我们日用的食粮；宽免我们的罪债，犹如我们也宽免得罪我们的人；不要让我们陷入诱惑}\BibleRef{路 11:2}，即陷入危险和罗网。既然祂命令他们祈祷\BibleQuote{不要让我们陷入诱惑}，祂便以实践这训诫来教导，说：\BibleQuote{圣父，若是可能，就让这杯离开我吧！}，以此教导所有圣人不要自投险境，不要冲入危险，而要等待危险临到，并显出一切勇气，只是不要自己冲前，或率先迎向恐怖。为何如此？为教导我们谦卑，也为我们免于虚浮光荣的指责。为此，这段经文也记载，祂说了这些话之后，\BibleQuote{便走去祈祷}；祈祷后，祂对门徒说：\BibleQuote{你们竟不能同我醒寤一个时辰吗？醒寤祈祷吧！免陷于诱惑。}\BibleRef{玛 26:39} 你看见祂不仅祈祷，也劝戒吗？祂说：\BibleQuote{心神固然切愿，但肉体却软弱。}\BibleRef{玛 26:41} 祂这样说，为虚怀他们的心，除去骄傲，教导他们节制，训练他们实践中庸。因此，祂愿意教导他们的祈祷，祂自己也献上，按人的方式说话，而非按祂的天主性（因为天主性是不受苦难的），而是按祂的人性。祂祈祷，是为教导我们祈祷，甚至求脱离困苦；但若不被允许，则要顺从天主看为好的。因此祂说：\BibleQuote{但不要照我的意愿，而要照你的意愿。}并非祂有一个意愿而圣父有另一个，而是为教导人，即使他们困苦战栗、危险临到、不愿脱离今生，也要将自己的意愿置于天主的意愿之后。正如保禄受教后也实际展现了这两个原则：他求主除去诱惑，说：\BibleQuote{关于这事，我曾三次求主使它脱离我}\BibleRef{格后 12:8}；但既然天主不愿除去，他便说：\BibleQuote{所以我甘愿在软弱、凌辱、迫害中喜乐。}\BibleRef{格后 12:10} 但也许我所说的不够清楚：因此我要说得更明白些。保禄遭遇了许多危险，并祈祷不要遭遇这些。然后他听到基督说：\BibleQuote{有我的恩宠为你够了，因为我的德能在软弱中才全显出来。}\BibleRef{格后 12:9} 他一见到天主的意愿是什么，今后便将自己的意愿服从于天主的意愿。因此，基督藉此祈祷教导了这两点：我们不应自投险境，而应祈祷不陷入其中；但若危险临到，我们应勇敢承受，并将自己的意愿置于天主的意愿之后。知道这些事，让我们祈祷永不陷入诱惑；但若陷入，让我们恳求天主赐予耐心和勇气，并尊崇祂的意愿高于我们自己的每一个意愿。这样，我们将平安度过今生，并获得将来的福乐：愿我们众人因我们的主耶稣基督的恩宠和慈爱，获得这一切，愿光荣、权能、尊崇归于祂，偕同圣父及圣神，从现在直到永远，世世无穷。阿们。\par

\SourceRecord{Translated by W.R.W. Stephens.；Nicene and Post-Nicene Fathers, First Series；Vol. 9.；Edited by Philip Schaff.；Buffalo, NY: Christian Literature Publishing Co.,；1889.；<http://www.newadvent.org/fathers/1910.htm>.}
